\begin{table}[!htbp]
\centering\small
\caption{Current study cohorts and configurations}\label{tab:1}
\begin{tabular}{@{}*{5}{>{\RaggedRight\arraybackslash}p{\dimexpr\linewidth/5-2\tabcolsep\relax}}@{}}
\toprule
Study & Cases & Assignments & Conditions & Scope \\
\midrule
Authorization core & 180 & 720 & Policy/\allowbreak{}full; auto/\allowbreak{}exact target & Pro; normal prompt \\
Nested stress regimes & 60 (nested) & 180 & Full; broad/\allowbreak{}underspecified/\allowbreak{}conflicting & Pro; normal prompt \\
Con\-tem\-po\-ra\-ne\-ous calibration & 120 & 480 & ND/\allowbreak{}policy/\allowbreak{}full/\allowbreak{}configured VIGIL & Pro; normal prompt \\
Synthetic controls & 4 scenarios & 16 & Policy/\allowbreak{}full & Pro + V4.1 Flash \\
\bottomrule
\end{tabular}
\par\smallskip\RaggedRight
\textit{Note.} The 900 authorization assignments use 180 cases representing 153 underlying tasks. All 120 calibration cases and 60 stress cases are nested within those 180 cases. The studies total 1,396 executions including 16 synthetic controls; the review reuses 240 records. Exact target is evaluator-derived oracle authorization. VIGIL uses generated constraints and deterministic hybrid auditing with secondary model verification disabled; its call cap is 16 versus 8 for the other calibration conditions.
\end{table}
